Surtout, on les mettait souvent à contribution pour donner un coup de main aux
pompiers en cas de feux des forêts et, bien sûr, ils exploitaient la forêt pour le compte
de l’Administration. Leur salaire était minime.
Lorsque les hameaux furent supprimés en 1965, plusieurs familles furent envoyées à
Lodève dans des logements HLM, d’autres à Mende dans un foyer Sonacotra, d’autres
en Provence.

A la fin Novembre 2006 j’ai pu glaner d’autres renseignements sur les Harkis de
Meyrueis et de Lozère auprès des deux Inspecteurs des Eaux et Forêts, successivement
en poste à Mende à cette époque : M. Paul Cabane et M. Marcel Gavalda..

C’est à \textbf{M. Cabane} qu’échut la responsabilité, dans les derniers mois de
1962, d’installer les harkis et leurs familles, parfois très nombreuses dans des
implantations choisies par la Direction des Eaux et Forêts en Lozère, dans le Gard ou
en Provence, à Meyrueis en ce qui nous concerne..

\begin{quote}
\textit{« Le grand problème fut de les loger, me dit-il, et les baraques ou les cabanes
étaient très légères. Ce ne devait être que des logements provisoires pour quelques
mois en attendant des installations définitives : en fait ils restèrent bien plus
longtemps. Seuls furent assez vite supprimés les hameaux qui étaient en altitude à
l’ombre et au froid.}

\textit{Lors de leur arrivée il n’y eut pas de problème pour leur installation. Et au début,
tout cela fut assez positif. Mais des agents d’encadrement ou des hommes de
l’extérieur, chevilles de syndicats les « chauffèrent de plus en plus » Ils travaillèrent
de moins en moins sur les chantiers forestiers, les pistes d’exploitation, les vacations
avec les pompiers, et ils en demandèrent de plus en plus. Et bien sûr on ne pouvait
approcher les mères ni leurs enfants »}
\end{quote}

\textbf{M. Gavalda} arrive à Mende en 1967. Donc trois ans après la mise en place des
hameaux forestiers.

\begin{quote}
\textit{« Je garde un souvenir ému des premiers contacts que j’eus avec eux. Je venais de
passer 15 ans au milieu des berbères du Haut Atlas au Maroc. Je connaissais
quelques mots de berbère et j’eus grand plaisir à retrouver une majorité de kabyles
auprès desquels je n’hésitais pas à revêtir djellaba ou gandoura.}

\textit{Il appartenait au service des Eaux et Forêts de choisir les chantiers, de définir les
travaux, les plantations, l’ouverture ou l’entretien de pistes forestières. Je disposais
d’une ligne budgétaire très encadrée pour les salaires, le matériel...}

\textit{Mais surtout ils devinrent « fantassins du feu », auxiliaires des pompiers, qui ne
disposaient que de bien peu de matériel : les camions modernes bien équipés
n’existaient pas, ni, à plus forte raison, les hélicoptères bombardiers d’eau ou les
Canadairs.}

\textit{Les harkis avaient seulement des pelles, des pioches, des battes pour frapper le feu
ou des haches pour faire des coupe-feux ou des contre-feux, et les incendies dans le
massif de l’Aigoual étaient fréquents et violents. A cette époque on n’avait pas
encore inventé les « tours de guet » et on ne pouvait surveiller les départs de feu
comme on peut le faire aujourd’hui avec les petits avions d’observation.}
\end{quote}
